\subsection{\texorpdfstring{复指数函数、\(\cos x\) 与 \(\sin x\) 的展开。}{复指数函数、cos x 与 sin x 的展开。}}

设 z 是任意复数，考虑实变量 t 的函数 \(\varphi(t) = e^{zt}\) ；我们有 \(\mathrm{D}^{n}\varphi(t) = z^{n}e^{zt}\) 与 \(e^{z} = \varphi(1)\) ；于是，用 \(\varphi(1)\) 在点 t = 0 处的 n 阶泰勒公式来表示它（II, p.~62），便得

\[
e ^ {z} = 1 + \frac {z}{1 !} + \frac {z ^ {2}}{2 !} + \dots + \frac {z ^ {n}}{n !} + z ^ {n + 1} \int_ {0} ^ {1} \frac {(1 - t) ^ {n}}{n !} e ^ {z t} d t\tag{6}
\]

当 \(z\) 为实数时这个公式与 (1) 等价。余项

\[
r _ {n} (z) = z ^ {n + 1} \int_ {0} ^ {1} \frac {(1 - t) ^ {n}}{n !} e ^ {z t} d t
\]

的绝对值可用中值不等式从上方控制；若 \(z = x + iy\) ，我们有 \(|e^{zt}| = e^{xt}\) ，于是 \(|e^{zt}| \leqslant 1\) （当 \(x \leqslant 0\) ），\(|e^{zt}| \leqslant e^x\) （当 \(x > 0\) ）；因此

\[
\left| r _ {n} (z) \right| \leqslant \frac {\left| z \right| ^ {n + 1}}{(n + 1) !} \text {   if   } x \leqslant 0\tag{7}
\]

\[
| r _ {n} (z) | \leqslant \frac {| z | ^ {n + 1} e ^ {x}}{(n + 1) !} \text {   if   } x > 0.\tag{8}
\]

与上面一样，我们得出

\[
e ^ {z} = \sum_ {n = 0} ^ {\infty} \frac {z ^ {n}}{n !}\tag{9}
\]

该级数在 \(\mathbf{C}\) 的每个紧子集上绝对且一致收敛。由 (6) 特别得出

\[
e ^ {i x} = 1 + \frac {i x}{1 !} + \frac {i ^ {2} x ^ {2}}{2 !} + \dots + \frac {i ^ {n} x ^ {n}}{n !} + i ^ {n + 1} \int_ {0} ^ {x} \frac {(x - t) ^ {n}}{n !} e ^ {i t} d t\tag{10}
\]

由此可得 \(\cos x\) 与 \(\sin x\) 的泰勒展开；取 (10) 的 \(2n + 1\) 阶实部便得

\[
\cos x = 1 - \frac {x ^ {2}}{2 !} + \dots + (- 1) ^ {n} \frac {x ^ {2 n}}{(2 n) !} + (- 1) ^ {n + 1} \int_ {0} ^ {x} \frac {(x - t) ^ {2 n + 1}}{(2 n + 1) !} \cos t d t\tag{11}
\]

其余项满足界

\[
\left| \int_ {0} ^ {x} \frac {(x - t) ^ {2 n + 1}}{(2 n + 1) !} \cos t d t \right| \leqslant \frac {| x | ^ {2 n + 2}}{(2 n + 2) !}.\tag{12}
\]

类似地，取 (10) 的 \(2n\) 阶虚部，我们得到

\[
\begin{array}{l} \sin x = x - \frac {x ^ {3}}{3 !} + \frac {x ^ {5}}{5 !} + \dots + (- 1) ^ {n - 1} \frac {x ^ {2 n - 1}}{(2 n - 1) !} \\ \qquad + (- 1) ^ {n} \int_ {0} ^ {x} \frac {(x - t) ^ {2 n}}{(2 n) !} \cos t d t \end{array}\tag{13}
\]

其余项满足界

\[
\left| \int_ {0} ^ {x} \frac {(x - t) ^ {2 n}}{(2 n) !} \cos t d t \right| \leqslant \frac {| x | ^ {2 n + 1}}{(2 n + 1) !}.\tag{14}
\]

此外，比较 (11) 中 \(2n + 1\) 阶与 \(2n + 3\) 阶的余项，有

\[
\int_ {0} ^ {x} \frac {(x - t) ^ {2 n + 3}}{(2 n + 3) !} \cos t d t = \frac {x ^ {2 n + 2}}{(2 n + 2) !} - \int_ {0} ^ {x} \frac {(x - t) ^ {2 n + 1}}{(2 n + 1) !} \cos t d t
\]

并把 (12) 考虑在内，便知不论 x 为何值，(11) 的余项都有与 \((-1)^{n+1}\) 相同的符号；同样可以证明 (13) 的余项与 \((-1)^{n}x\) 同号。特别地，在 (11) 中取 n = 0 与 n = 1，在 (13) 中取 n = 1 与 n = 2，就得到不等式

\[
1 - \frac {x ^ {2}}{2} \leqslant \cos x \leqslant 1 \quad \text { for   all } x\tag{15}
\]

\[
x - \frac {x ^ {3}}{6} \leqslant \sin x \leqslant x \quad \text { for   all } x \geqslant 0.\tag{16}
\]

最后，在 (9) 中令 z = ix 便得

\[
\cos x = \sum_ {n = 0} ^ {\infty} (- 1) ^ {n} \frac {x ^ {2 n}}{(2 n) !}\tag{17}
\]

\[
\sin x = \sum_ {n = 0} ^ {\infty} (- 1) ^ {n} \frac {x ^ {2 n + 1}}{(2 n + 1) !}\tag{18}
\]

这些级数在每个紧区间上绝对且一致收敛。

此外，显然公式 (17) 与 (18) 对每个复数 x 仍然成立，右端的级数在 C 的每个紧子集上绝对且一致收敛。特别地，对每个 x（实的或复的）

\[
\begin{array}{l} \cosh x = \sum_ {n = 0} ^ {\infty} \frac {x ^ {2 n}}{(2 n) !} \\ \sinh x = \sum_ {n = 0} ^ {\infty} \frac {x ^ {2 n + 1}}{(2 n + 1) !}. \end{array}
\]

